% Round 21 / Hodge branch. New candidate; independent review required.
\section{A mixed-degeneration global family with complete heights}
\label{r21mg:section}

\paragraph{Precise contribution.}
The local MX21 theorem can replace the distinguished-place input in
the frozen SG18 coefficient construction.  We carry out that
replacement with literal integral coefficients and check every finite
place.  The resulting family has the full homologically trivial Chow
groups as its height domain after every number-field extension.
It is not the SG18 family with a different model: at the distinguished
places its fourth monodromy is zero, whereas the fourth monodromy of
every SG18 member has rank $179$.  It also has a nonzero weight-two
piece, so it has no potential good reduction there.

This does not classify overlap with every earlier abstract motivic
criterion, and it does not prove G for every member of the ambient
linear system or for arbitrary varieties.  No historical priority is
claimed for the classical arithmetic-pencil or semistable-resolution
methods used in the construction.

\subsection{The exact frozen global input and its check}

Use the global data of frozen SG18:
\[
 E_0/\mathbf Q:\ y^2=x^3-x,\qquad
 F_0=\mathbf Q(E_0[3]),\qquad
 \mathcal S=\mathcal E\times_{\mathcal O_{F_0}}\mathcal E,
 \qquad
 \mathcal A_0=\mathcal O_{\mathcal E}(3[0])
                     \boxtimes\mathcal O_{\mathcal E}(3[0]).
 \tag{MG21.1}\label{r21mg:surface}
\]
Here $F_0$ contains the coordinates of \emph{all} geometric
$3$-torsion points.  SG18 proves that $F_0/\mathbf Q$ is finite
Galois and unramified at $5$, that $E_0$ has good reduction at
every finite place of $F_0$, and that $\mathcal E$ is a smooth
projective elliptic scheme over the entire number ring.

We recall the hypotheses and degree checks in that proof to make
the reuse precise.  The discriminant of the displayed equation is
$64$, so only places above $2$ require a new argument.  Complex
multiplication over $\mathbf Q(i)$ gives potential good reduction
there by Serre--Tate, Theorem 6(a).  Thus the $3$-adic inertia image
at $2$ is finite.  Over $F_0$ it is contained in
$1+3M_2(\mathbf Z_3)$, since all $3$-torsion is rational.
This congruence group has no nontrivial finite-order elements:
for $1+3^aB\ne1$ the binomial expansion of a prime-order power
has a unique lowest-valuation term, including at the prime $3$.
The finite inertia image is therefore trivial.  The prime-to-residue
Neron--Ogg--Shafarevich criterion, Serre--Tate Theorem 1, supplies
good reduction at $2$.  Good reduction already holds at every other
prime, including $3$.  At $5$ the original elliptic model has finite
etale $3$-torsion, so inertia on the full torsion field is trivial.
The local proper smooth elliptic models glue by the Neron property;
the identity divisor of relative degree three gives projectivity.
This does not assert that the same Weierstrass equation is smooth
over the whole number ring.

Set
\[
  \mathcal Y=\mathcal S\times\mathbf P^3_{\mathcal O_{F_0}},
  \qquad \mathcal L_0=\mathcal A_0\boxtimes\mathcal O(1),
  \qquad \mathcal A=\mathcal A_0^5.
 \tag{MG21.2}\label{r21mg:ambient}
\]
All these schemes are smooth projective with geometrically integral
fibers at every finite place.  In particular
$\mathcal A\boxtimes\mathcal O(5)=\mathcal L_0^5$, exactly the
line bundle in the frozen SG18 global coefficient space.

The retained SG18 source is
\texttt{research\_input/round18/sectors/}
\texttt{abelian\_surface\_ambient\_complete\_heights.tex},
SHA-256
\texttt{14fb50ae6121f72918f18850c4b8e8c17ac7289ca8305c9205cb7497751f3905}.
Its Lemmas \texttt{r18sg:global-good} and
\texttt{r18sg:coefficients}, including their native-field scope,
are the exact global inputs here.  The original Serre--Tate primary
source is
\url{https://wstein.org/papers/bib/Serre-Tate-Good_Reduction_of_Abelian_Varieties.pdf},
printed pages 493 and 503.

\begin{lemma}[Actual free coefficients for the mixed product]
\label{r21mg:coefficients}
There is a fixed relative projective embedding
$\mathcal Y\hookrightarrow\mathbf P^N_{\mathcal O_{F_0}}$ with
$\mathcal O(1)|_{\mathcal Y}=\mathcal L_0$ such that the finite free
module $V_5$ of degree-five ambient polynomials surjects onto
$H^0(\mathcal Y,\mathcal L_0^5)$ and onto its section space on
every field fiber.  The same remains true after the finite field
extension used below.  Consequently every mixed product
\[
       P=f_0\prod_{j=1}^5l_j,
       \quad f_0\in H^0(\mathcal S,\mathcal A_0^5),
       \quad l_j\in H^0(\mathbf P^3,\mathcal O(1)),
 \tag{MG21.3}\label{r21mg:P}
\]
has an actual degree-five polynomial representative in $V_5$.
\end{lemma}

\begin{proof}
This is SG18's full-section coefficient lemma, with the mixed
product interpreted in its correct line bundle.  For clarity, on
each elliptic fiber the degree-three embedding is a plane cubic;
$H^1(\mathbf P^2,\mathcal O(t))=0$ makes its coordinate polynomials
surject onto every positive tensor degree.  Combining the two
elliptic embeddings and $\mathbf P^3$ by Segre gives the same
surjectivity for $\mathcal L_0^d$ in every field fiber.  Positive
elliptic line bundles have no higher cohomology.  Kunneth and
cohomology and base change make the section module locally free,
and Nakayama proves the global surjectivity from a chosen finite
free ambient coordinate module.  No freeness of the section module
itself over a Dedekind ring is required.  These arguments commute
with the finite flat number-ring base change.

The two factors in \eqref{r21mg:P} have bidegrees $(5,0)$ and
$(0,5)$ relative to $\mathcal A_0$ and $\mathcal O(1)$; their product
therefore belongs to $H^0(\mathcal L_0^5)$.  Surjectivity gives the
asserted polynomial representative.  We choose one such representative
before imposing any coefficient congruence.  The individual factors
are not incorrectly regarded as degree-one Segre forms.
\end{proof}

\subsection{The complete global family}

\begin{theorem}[Mixed global complete-height family, MG21]
\label{r21mg:main}
There is one fixed finite extension $F/F_0$, unramified at all places
over $5$, and a Zariski-dense family of smooth geometrically integral
fourfolds
\[
                              X\in|\mathcal L_0^5|_{Y_F}
 \tag{MG21.4}\label{r21mg:family}
\]
with the following properties.
\begin{enumerate}
\item For every finite extension $F'/F$, every finite place $w$ of
$F'$, every codimension, and every specified geometrically
homologically trivial rational cycle on $X_{F'_w}$, strong G holds
with exact support, all compact ordinary actual tests, and the given
absolute value.
\item For every such $F'/F$, the frozen PH16 pairings are defined on
the complete complementary geometric homological Chow groups:
\[
 \operatorname{CH}^r(X_{F'})_{\mathbf Q,\mathrm{geom\,hom}}
   \times\operatorname{CH}^{5-r}(X_{F'})_{\mathbf Q,\mathrm{geom\,hom}}
          \longrightarrow\mathbf R,\qquad 1\le r\le4.
 \tag{MG21.5}\label{r21mg:height}
\]
GA17 supplies regular-model admissible Chow lifts on these full
domains.
\item At every place of $F$ over $5$, the actual MX21 model has at
most four components through a point, the fourth monodromy operator
on $H^4$ is zero, and its weight-two graded piece contains a
nonzero Frobenius-stable direct summand of dimension $8$.
The generic fiber has no potential good reduction at these places.
\item After a common finite extension of the corresponding local
fields at $5$, no member is geometrically isomorphic to a member of
the frozen explicit SG18 family, whose fourth monodromy has rank
$179$.  This distinguishes these generic fibers, not merely their
chosen models.
\end{enumerate}
The density is over the one fixed $F$ and in the full section space.
No assertion that every member of that linear system has G is made.
\end{theorem}

\begin{proof}
\emph{1. Fix the reference and all residue choices before the field.}
Choose an integral polynomial $H_0\in V_5$ whose generic divisor
on $Y_{F_0}$ is smooth.  The very ample degree-five system has a
nonempty smooth Bertini open, and an integral coefficient lattice
is Zariski dense, so such a choice exists.
Let $S_0$ be the finite set of rational primes other than $5$ lying
below the places where that divisor is not smooth, together with
$2$.  Include a place where its entire defining section vanishes
on reduction in this set as well.

At \emph{every} place $v$ of $F_0$ over every prime in $S_0$, choose
a smooth degree-five special section $G_v$ after a finite residue
extension.  The smooth open is geometrically nonempty.  It is
essential to make this choice even at places over the same rational
prime where $H_0$ was already smooth.

At every $v\mid5$, choose instead, after a finite residue extension,
a smooth geometrically integral curve cut by
$f_{0,v}\in H^0(S_v,\mathcal A_{0,v}^5)$, five hyperplanes
$l_{j,v}$ of $\mathbf P^3$ in general position, and
$Q_v\in H^0(Y_v,\mathcal L_{0,v}^5)$ transverse to every mixed
stratum, with the geometric integrality and finite-etale conditions
of MX21.  The line bundles are very ample, so simultaneous Bertini
on the finite list of strata gives a nonempty open over the
algebraic closure and hence a point over a finite residue extension.

Choose an integer $N$ divisible by all the finitely many residue
extension degrees just chosen.  At each of the selected maximal
ideals of $\mathcal O_{F_0}$ choose a monic irreducible polynomial
of degree $N$ over its residue field.  Coefficientwise CRT gives one
monic $g\in\mathcal O_{F_0}[T]$ with all these reductions.  It is
irreducible over $F_0$: a monic factor would have integral
coefficients in the integrally closed number ring and would give
a nontrivial factor of one of the irreducible reductions.
Put $F=F_0[T]/(g)$.  At each selected place this extension is
unramified and inert of degree $N$, since the reduction of $g$
is separable and irreducible.  The corresponding finite etale local
algebra is already the unramified integer ring.  Thus all selected
residue fields are available over this fixed $F$, and $F/\mathbf Q$
is unramified at $5$.  The everywhere-good surface base changes to
an everywhere-good surface over $\mathcal O_F$.

\emph{2. Lift actual mixed factors and impose literal congruences.}
The section module of $\mathcal A_0^5$ on the surface commutes with
reduction, by the positive-degree elliptic cohomology calculation.
CRT in this module lifts the finitely many $f_{0,v}$ at $v\mid5$
to one global section $f_0$.  The hyperplane coefficients lift in
the free module $\mathcal O_F^4$, and $Q_v$ lifts in the ambient
free module $V_5$.  Form the actual mixed section
$P=f_0\prod_{j=1}^5l_j$ and choose a polynomial representative
$P\in V_5$ using Lemma~\ref{r21mg:coefficients}.  Hence all
congruences below are literal congruences of ambient polynomials.
Use the same coefficient surjectivity to represent the chosen
$G_v$ at every place over $S_0$.

Put $B=\prod_{p\in S_0}p$ and $M=25B$.  Choose $H_1\in V_5$
subject to
\[
 H_1\equiv P+5Q\pmod{25\mathcal O_F},\qquad
 H_1\equiv G_v\pmod{\mathfrak p_v}\quad(v\text{ above }S_0).
 \tag{MG21.6}\label{r21mg:H1}
\]
These conditions are pairwise compatible.  At $5$ the equality
$25\mathcal O_F=\prod_{v\mid5}\mathfrak p_v^2$ holds because
the field is unramified there.  Their solution set is a nonempty
translate of a full integral coefficient lattice.

Require also that the pencil $[uH_0+vH_1]$ have smooth base locus
and only transverse one-node singular geometric members, and that
$H_1$ itself be smooth.  The fixed-smooth-member Lefschetz open
proved in frozen CG18 and used in SG18 is nonempty on this actual
smooth fivefold.  Its proof needs only first and second jets at one
point and independent first jets at two distinct points.  Degree-five
polynomials in the given projective embedding supply all these jets.
The marked-singularity incidence gives a discriminant hypersurface
with a dense one-node open; projection from the fixed smooth
$[H_0]$ is finite and generically etale in characteristic zero.
Avoiding its bad and branch loci gives exactly the stated pencil
condition.  No step uses how the residue target $P$ factors.
The dense congruence lattice therefore meets this open densely.

\emph{3. Spread each chosen pencil and select infinitely many parameters.}
After $H_1$ is fixed, spread the pencil and enlarge a finite set of
places so that outside it the base locus is smooth and every
geometric fiber is smooth or has exactly one ordinary quadratic
point.  This is the same finite-presentation argument as SG18:
the finite critical-value scheme, its separation of distinct values,
the nonzero Hessians, and the smoothness off it spread after
inverting finitely many primes.  The two coefficient sections also
remain independent there, so no parameter gives a zero defining
section.  Include all places above $S_0\cup\{5\}$ and all
characteristic-two places in the exceptional set.
Let $T$ be the finite set of underlying rational primes.  It may
depend on $H_1$, but neither $F$ nor the earlier lattice does.

For every integer $a$ satisfying
\[
 a\equiv1\pmod{25},\qquad a\equiv1\pmod B,\qquad
 a\equiv0\pmod p\quad(p\in T\setminus(S_0\cup\{5\})),
 \tag{MG21.7}\label{r21mg:a}
\]
take the divisor
\[
                              X_a=V(MH_0+aH_1)\subset Y_F.
 \tag{MG21.8}\label{r21mg:Xa}
\]
The CRT moduli are coprime after repetitions are removed, so there
are infinitely many such integers.  Remove the finitely many which
give a singular generic member.  All remaining members are smooth.

\emph{4. Check the actual local model at every finite place.}
At $v\mid5$, write the literal polynomial equality
$H_1=P+5Q+25R_1$.  The integer $a$ is a unit there.  Dividing the
local equation by $a$ gives
\[
 P+5Q_a=0,\qquad
 Q_a=Q+5R_1+5B a^{-1}H_0,
                       \qquad\bar Q_a=\bar Q.
 \tag{MG21.9}\label{r21mg:localmixed}
\]
Here $5$ is a native uniformizer.  This is exactly the MX21
mixed model with $\mathcal A=\mathcal A_0^5$ and $d=5$,
including the previously chosen actual transverse strata.

At every place above $S_0$, the integer $M$ reduces to zero and
$a$ to one.  The special defining section is the chosen smooth
$G_v$.  At any exceptional place outside $S_0\cup\{5\}$, $M$
is a unit and $a$ reduces to zero, so the special divisor is the
smooth reference $H_0$.  This reference is smooth there because
all of its bad rational primes were already put into $S_0$.
At every other place the spread pencil has a smooth or one-node
special divisor.  Thus there is no missing finite place.

Every defining special section just described is nonzero.  In the
smooth ambient with integral fibers it is a nonzerodivisor; the
relative Cartier divisor is flat over the trait.  The generic
fiber is geometrically integral, and flatness plus the Cartier
Cohen--Macaulay property makes the total closure integral.  The
required local projective models are therefore actual flat
integral models, not just descriptions of individual equations.

For clarity, smooth and one-node geometric fibers are integral.
On every field fiber,
\[
 H^1(S\times\mathbf P^3,
                \mathcal A_0^{-5}\boxtimes\mathcal O(-5))=0:
 \tag{MG21.10}\label{r21mg:connected}
\]
the two Kunneth terms contain respectively
$H^0(\mathbf P^3,\mathcal O(-5))=0$ and
$H^1(\mathbf P^3,\mathcal O(-5))=0$.
The divisor sequence gives $H^0(X,\mathcal O_X)=\bar k$ after
geometric base change.  Such a Cartier fourfold is $S_2$, and
an isolated-node fiber is also $R_1$, hence normal.  Connectedness
and normality give geometric integrality.  The same reasoning
applies to the smooth generic members.

\emph{5. Activate every complete local cycle domain, then the heights.}
At every place above $5$, MX21 proves full G and its all-finite-
extension version for every new cycle over that extension.
At every other place the literal model above meets the frozen
good-reduction or ON17 ordinary-quadratic theorem, including its
all-finite-extension statement.  Characteristic-two places were
forced smooth.  For an arbitrary finite $F'/F$ and every $w$ over
these places, apply the corresponding local theorem to the entire
new Chow group on $X_{F'_w}$.  No such group is replaced by cycles
pulled back from $F$.

The frozen PH16 hypotheses now hold at all finite places in every
complementary degree.  Together with its established archimedean
input, PH16 yields \eqref{r21mg:height} on the full original global
homological Chow groups over $F'$.  GA17 supplies the admissible
regular-model lifts on this proved domain.  All actual local
identities use their exact specified supports, all compact actual
tests, the full auxiliary field degrees, and the fixed local
absolute values.

\emph{6. Prove density over the one fixed number field.}
The choices of $F$ and of the residue coefficient lattice precede
the variation of $H_1$.  The allowed $H_1$ in that lattice and in
the fixed-member open are Zariski dense in $V_{5,F}$.
For each such direction, infinitely many integers $a$ in
\eqref{r21mg:a} survive the finite exclusion.  A polynomial
vanishing on all constructed coefficient vectors $MH_0+aH_1$
therefore vanishes on the whole affine line in each such direction.
In particular it vanishes at $MH_0+H_1$.  Those latter points
are dense, so the polynomial is zero.  The surjection of
Lemma~\ref{r21mg:coefficients} makes the images dense in the full
section space, and hence in the displayed projective linear system.

\emph{7. Check the genuine generic-fiber distinction.}
The sign in \eqref{r21mg:surface} matters: this is $y^2=x^3-x$,
not the plus-sign illustration in MX21.  At $5$ it has eight points,
Frobenius trace $-2$, and eigenvalues $\alpha=-1+2i$, $\beta=-1-2i$
on $H^1$.  Its square surface has four independent divisor classes
from the factors, the diagonal, and the order-four automorphism
graph; their intersection matrix has determinant $-4$.
The other two $H^2$ eigenvalues are
\[
                   \alpha^2=-3-4i,\qquad \beta^2=-3+4i.
 \tag{MG21.11}\label{r21mg:T}
\]
Their quotients by $5$ are not roots of unity: they lie in
$\mathbf Q(i)$ and are not $1,-1,i,-i$.  A geometric divisor is
defined over some finite residue extension and has a Tate
eigenvalue there.  Thus these two directions remain outside the
divisor subspace after every residue extension, while the four
displayed independent divisors fill the other four-dimensional
space.  Hence $\dim T=2$ at every completion above $5$ used here.

MX21 with $d=5$ gives a direct summand of dimension
$2\binom43=8$ in the weight-two part of $H^4$ and therefore excludes
potential good reduction.  Its actual model has no fivefold
stratum.  The frozen NM17 crossing-depth bound gives $N^4=0$.
Equivalently, the weight spectral sequence has no column outside
$-3\le a\le3$, and $N$ lowers the abutment weight filtration by
two; its fourth power is zero on $H^4$.

Frozen SG18 instead proves $\operatorname{rank}N^4=179$ at every
place above $5$ for every member of its explicitly constructed
family.  Restriction to a finite valued extension multiplies a
normalized logarithmic monodromy by a nonzero ramification scalar,
so zero versus nonzero fourth power persists.  A geometric
isomorphism over an algebraic closure is defined over a finite
extension and would identify the corresponding cohomological
Galois representations after a common finite extension.  The two
monodromy assertions contradict that possibility.  This proves
the stated strict distinction from the named frozen SG18 family.
\end{proof}

\subsection{Boundaries of the theorem}

The new global family has the same complete cycle-group scope as
SG18 but a different certified local degeneration type and different
intrinsic monodromy at $5$.  In particular the calculation proves
more than another model of the same SG18 generic fibers.  It does
not assert disjointness from every earlier abstract correspondence,
domination or motivic sufficient criterion; those are larger
collections without a complete isomorphism classification.

Weak Lefschetz still gives $b_1(X)=4$ and $b_2(X)=7$, since the
ambient and total line bundle are unchanged.  Consequently these
are not smooth projective-space complete-intersection fourfolds,
nor members of the frozen projective-space hypersurface families.
The comparison with SG18 uses $N^4$, not these shared Betti numbers.

The extension from local G to complete heights was made only after
checking every finite place in the construction.  Neither a local
dense family alone nor the existence of a positive model at one
place would have sufficed.  General G for arbitrary varieties and
the arbitrary-variety complete unconditional BB height problem
remain \textbf{OPEN}.

\paragraph{Dependency ledger.}
The genuinely new local input is MX21.  The preserved frozen
arithmetic inputs are SG18's everywhere-good surface and free
coefficient lemma, the CG18/SC18 fixed-smooth-member Lefschetz
pencil proof, the finite-place spreading and CRT mechanism written
out above, ON17 and the frozen good-reduction all-cycle theorems,
RI18 as used inside MX21, PH16, and GA17.  The distinction uses
NM17's crossing-depth bound and SG18's checked fourth-monodromy
rank.  The $5$-adic transcendental eigenspaces are recomputed above;
no identification of the two different elliptic equations is
assumed.
